\subsubsection{Uncertainty-Aware Traversability Cost Map}\label{sec:costmap}

% Motivation
Off-road environments are fraught with bumps and obstacles of varying shapes, despite being in the same terrain class. For safe and effective navigation of a vehicle in off-road terrain, estimating the nuanced traverse cost of the terrain is necessary. A path planner can optimize a trajectory that minimizes disturbances during navigation with the predicted cost. Therefore, we generate dense and continuous traverse cost maps in bird's eye view~(BEV) from a single LiDAR point cloud.

% data generation
The z-axis linear acceleration measured by an Inertial Measurement Unit~(IMU) is used to define the traversability that can stabilize navigation, as this component effectively communicates the terrain's bumpiness~\cite{Bekhti_verticala}. In addition, because our controller employs a vehicle dynamics model that is ignorant of vertical motions~\cite{kim_smooth_2022, kim_physics_2022}, the definition of traversability based on vertical acceleration can be helpful in preventing control performance losses. Motivated by recent work~\cite{castro2022does}, we define traverse cost using the spectral power of z-acceleration produced by a continuous wavelet transformation. 

\begin{figure}[t]
\centering
\includegraphics[width=1.0\linewidth]{Figures_Final/pipeline_compressed.pdf}
\caption{Overview pipeline of the traversability cost prediction network.}
\label{fig:pipeline}
\vspace*{-0.25in}
\end{figure}

% base network structure
The network is trained to produce a dense traverse cost map using a sparse single-sweep LiDAR point cloud in BEV, as shown in Fig.~\ref{fig:pipeline}. Following PointPillars~\cite{lang2019pointpillars}, each point is discretized into sparse voxels with a grid size of 0.2m. Using PointNet~\cite{qi2017pointnet}, each voxel is converted into sparse pillar features. Then, a U-Net~\cite{unet} structured network, which has an encoder-decoder architecture with skip connections, is employed to generate a dense pillar feature map. It progressively reduces the spatial size of features and captures higher-level semantic information, while the decoder upsamples feature maps to recover spatial information. Then, the dense pillar features are then concatenated with parameterized velocity to produce velocity-conditioned cost maps. The velocity is parameterized with Fourier feature mapping~\cite{tancik2020fourier} to incorporate vehicle's velocity into the costmap prediction~\cite{castro2022does}.


Finally, the MLP head predicts the mean ${\mu}_i$ and standard deviation ${\sigma}_i$ of the traversability for each pillar $i$. The network is trained to minimize the Gaussian log-likelihood:
\begin{equation}\label{gaussiannll}
    \mathcal{L}^{\text{traverse}}\left(\tau, \bm{\theta}\right) = \frac{1}{2}\left(\log({\sigma}_i) +  \frac{({\mu}_i - c_i)^2}{{\sigma}_i} \right),
\end{equation}where $\tau$, $c_i$, and $\bm{\theta}$ represent a trajectory, the cost of the the trajectory that is associated with the pillar $i$, and model parameter, respectively.

\subsubsection{Online Adaptation for Minimizing Aleatoric Uncertainty}\label{sec:meta}

\begin{algorithm}[h]
\footnotesize
\SetKwInOut{Input}{Given}
\Input{$\bm{\theta}$:
$\mathcal{D}$: Interaction data from various environments\;
$M, K$: Number of past and future timesteps\;
$N$: Number of sampled trajectories\;
$N_A$: Number of the inner gradient steps\;
$\alpha, \beta$: Learning rate for inner and outer gradient steps\;
Randomly initialize $\bm{\theta}$\;
\For{$i \leftarrow 0$ \KwTo \text{maximum iterations}}{
    \For{$j \leftarrow 0$ \KwTo $N-1$}{
        Sample $\tau(t-M, t)$, $\tau(t, t+K)$ $\sim \mathcal{D}$\;
        Self-Label $\tau(t-M, t)$ and $\tau(t, t+K)$\;
        $\bm{\theta}' \leftarrow \bm{\theta}$\;
        % start adaptation
        \For{$k \leftarrow 0$ \KwTo $N_A-1$}{
            $\bm{\theta}' \leftarrow \bm{\theta}' - \alpha \nabla_{\bm{\theta'}} \mathcal{L}^{\text{traverse}}\left(\tau(t-M, t), \bm{\theta'}\right)$\;
        }
        $\mathcal{L}_j \leftarrow \mathcal{L}^{\text{traverse}}\left(\tau(t, t+K), \bm{\theta'}\right)$
    }
    $\bm{\theta} \leftarrow  \bm{\theta} - \beta \nabla_{\bm{\theta}} \frac{1}{N} \sum\limits_{j=1}^N \mathcal{L}_j$
}
}
\caption{Meta Learning of Traverse Cost \label{Algorithm:meta}} 
\end{algorithm} 


% Uncertainty and why
A significant degree of aleatoric uncertainty occurs when learning traverse cost from a single-sweep LiDAR point cloud. In real-world off-road environments, there are a variety of terrain types, and the terrain and vehicle properties that influence traversability (e.g., deformability, friction, and roughness) differ according to a number of factors. However, such subtle differences cannot be precisely captured by a LiDAR point cloud. The ground-truth traverse cost of terrain captured in different environments would vary considerably, while the geometric characteristics of terrain would be comparable. Learning a global traversability model using a large dataset~$\mathcal{D}$ of multiple environments leads to high aleatoric uncertainty in estimation.

% To resolve -> MAML.(Motivation)
To resolve this issue, we propose a model that can effectively learn a global traversability model that is capable of quickly adapting to a new environment based on its recent experiences. We use Model-Agnostic Meta-Learning~(MAML)~\cite{finn2017model} for online adaptation of the traversability model. It aims to find the initial parameters of the network so that adaptation with a few gradient descent steps from this initialization leads to effective generalization to current circumstances. By doing so, the model can incorporate driving data collected in various environments for learning traverse cost from geometric properties without confusion.



% MAML Explanation
While terrain properties vary significantly in different environments, we assume that the environment is locally consistent. Consequently, each segment of a trajectory is regarded as a separate \textit{task}, denoted as $\tau$. The traversability model is trained to adapt using the meta-train data of $M$ past timesteps, $\tau\left(t-M,t\right)$, to predict the traverse cost of the next $K$ timesteps, $\tau\left(t,t+K\right)$, as follows: 
\begin{equation}\label{eq:meta}
    \begin{aligned}
    \argmin_{\bm{\theta}} \hspace{10pt} & \mathbb{E}_{\tau\left(t-M, t+K\right) \sim \mathcal{D}}  \big[ \mathcal{L}^{\text{traverse}}\left(\tau(t, t+K), \bm{\theta}'\right)\big] \\   
    & \text{s.t.:} \hspace{10pt} \bm{\theta}' = \bm{\theta} - \alpha \nabla_{\bm{\theta}} \mathcal{L}^{\text{traverse}}\left(\tau(t-M, t), \bm{\theta}\right).
    \end{aligned}
\end{equation} The past $M$ timesteps provide insight on how to adapt the model to precisely predict future trajectories' traverse costs. As illustrated in Fig.\ref{fig:concept}, the model can online-adapt during the deployment phase using automatically generated meta-train data in a self-supervised manner. Algorithm~\ref{Algorithm:meta} describes the meta-learning-based training procedure for the global traversability model. The trained global model is applicable in a variety of environments and is even able to adjust to unseen terrain. 
% MAML Algorithm



\begin{table}[b]
\vspace*{-0.1in}
\centering
\renewcommand {\arraystretch}{1.1}
\caption{Validation error of traversability estimation models for real-world driving data. Our method shows a significant margin compared to the baseline.}
\footnotesize
\label{tab:quantitative}
\resizebox{1.0\linewidth}{!}{%
    \begin{tabular}{cccccc}
        \toprule
        & \bf{Unpaved} & \bf{Rough} & \bf{Grassland} & \bf{Profiled Road} & \bf{Simulation} \\
        \midrule \midrule
        
        \multicolumn{1}{c|}{\textit{Baseline}}      & 0.1222 & 0.0196 & 0.2460 & 0.2039 & 0.7263 \\
        \multicolumn{1}{c|}{\textit{w.o. Adaptation}} & 0.0713 & 0.0171 & 0.1961 & 0.1907  & 0.6668 \\
        \multicolumn{1}{c|}{\textit{w. Adaptation}} & \textbf{0.0114} & \textbf{0.0015} & \textbf{0.1767} & \textbf{0.1523} & \textbf{0.5725} \\
\end{tabular}%
}
% \vspace*{-0.2in}
\end{table}